\section{Problem Summary}
Start from one positive integer \(n\). If the current value is even, divide it by \(2\); otherwise replace it with \(3n + 1\). Print every value in this chain until the sequence reaches \(1\).

\section{Editorial}
There is no optimization to discover here: the task is simply to simulate the process exactly as described.

Keep printing the current value, then apply one step of the rule. The loop stops once the current value becomes \(1\), and then \(1\) should also appear in the output.

The only practical detail is the numeric type. Even when the input fits comfortably in a 32-bit integer, the intermediate value \(3n + 1\) can grow past that range, so a 64-bit integer is the safe default.

\section{Pseudocode}
Read \(n\).

While \(n\) is not \(1\):
\begin{itemize}
\item print \(n\)
\item if \(n\) is even, replace it by \(n / 2\)
\item otherwise, replace it by \(3n + 1\)
\end{itemize}

Print \(1\).

\section{Complexity Analysis}
\begin{itemize}
\item Time: \(O(k)\), where \(k\) is the length of the produced sequence
\item Memory: \(O(1)\)
\end{itemize}

\section{Notes}
Use a 64-bit integer type so the odd-step update does not overflow.
